\chapter{Trigonométrie du triangle}
\label{workbook:chapter:15}
\section{Triangles rectangles}
\label{workbook:section:15.01}
\begin{exerciseblock}{Triangles rectangles}
\item Dans un triangle rectangle, l'hypoténuse vaut $13$ et un côté vaut $5$. Trouver l'autre côté et les deux angles aigus.
\item À $40$ m d'un bâtiment, l'angle d'élévation du sommet vaut $52^\circ$. Estimer la hauteur, en supposant l'instrument au sol.
\item Choisir entre sinus, cosinus et tangente pour trois situations : opposé/hypoténuse; adjacent/hypoténuse; opposé/adjacent.
\end{exerciseblock}
\section{Vue d'ensemble des triangles quelconques}
\label{workbook:section:15.02}
\begin{exerciseblock}{Vue d'ensemble des triangles quelconques}
\item Pour un triangle donné par $a=7$, $b=9$, $C=60^\circ$, choisir la loi appropriée et calculer $c$.
\item Pour $A=35^\circ$, $B=70^\circ$, $a=8$, calculer $C$, puis $b$ et $c$.
\item Classer les données SSS, SAS, ASA/AAS et SSA selon la loi à utiliser et le risque d'ambiguïté.
\end{exerciseblock}
\section{Choisir la formule à partir de la figure}
\label{workbook:section:15.03}
\begin{exerciseblock}{Choisir la formule à partir de la figure}
\item Pour quatre schémas correspondant à SSS, SAS, ASA et SSA, écrire seulement la première équation pertinente avant tout calcul.
\item Un triangle a $a=10$, $b=6$, $A=30^\circ$. Déterminer le nombre possible de triangles à l'aide de la hauteur critique.
\item Créer une liste de contrôles : somme des angles, côté opposé au plus grand angle, inégalité triangulaire.
\end{exerciseblock}
\section{Origine géométrique des lois du sinus et du cosinus}
\label{workbook:section:15.04}
\begin{exerciseblock}{Origine géométrique des lois du sinus et du cosinus}
\item Tracer une hauteur dans un triangle quelconque et dériver $a/\sin A=b/\sin B$.
\item À partir d'une projection, dériver $c^2=a^2+b^2-2ab\cos C$.
\item Montrer que la loi des cosinus devient Pythagore lorsque $C=90^\circ$.
\end{exerciseblock}
\section{Résoudre et valider un triangle}
\label{workbook:section:15.05}
\begin{exerciseblock}{Résoudre et valider un triangle}
\item Résoudre complètement un triangle avec $a=12$, $b=15$, $C=70^\circ$.
\item Analyser le cas $A=40^\circ$, $a=7$, $b=10$ : zéro, un ou deux triangles? Calculer les solutions possibles.
\item Après calcul, expliquer comment vérifier la cohérence du triangle sans refaire toute la résolution.
\end{exerciseblock}
\section*{Synthèse du chapitre}
\begin{synthesisbox}
\item Mesurer une largeur inaccessible : une base $AB=120$ m, angles $A=48^\circ$, $B=67^\circ$. Calculer les deux autres côtés et la hauteur sur $AB$.
\item Un triangle a côtés $5$, $7$, $10$. Calculer ses angles, son aire de deux façons et vérifier l'ordre des angles.
\item Construire un exemple numérique de cas SSA avec deux solutions, un avec une solution et un sans solution; justifier par la hauteur critique.
\end{synthesisbox}
